\documentclass[10pt,a4paper]{article}
\usepackage[margin=18mm,bottom=20mm]{geometry}
\usepackage{fontspec,xeCJK,amsmath,amssymb,xcolor,fancyhdr}
\setmainfont{Arial Unicode.ttf}[Path=/Library/Fonts/]
\setCJKmainfont{Arial Unicode.ttf}[Path=/Library/Fonts/]
\definecolor{ink}{HTML}{183852}
\definecolor{teal}{HTML}{146D80}
\setlength{\parindent}{0pt}
\setlength{\parskip}{5pt}
\linespread{1.16}
\setlength{\abovedisplayskip}{6pt}
\setlength{\belowdisplayskip}{6pt}
\pagestyle{fancy}\fancyhf{}
\fancyfoot[L]{\footnotesize RoomScout / 几何数据与数学模型 / LaTeX 版}
\fancyfoot[R]{\thepage}
\renewcommand{\headrulewidth}{0pt}
\renewcommand{\footrulewidth}{0.3pt}
\setlength{\emergencystretch}{3em}
\begin{document}

{\color{ink}\LARGE RoomScout\par 几何数据与数学模型说明\par}\vspace{5pt}
\par\vspace{5pt}{\color{teal}\large 源码对应版 · 2026-09-08\par}
本文把场景的数据结构、几何验证算法和修复选择机制对应到当前项目源码，说明每个模块接收什么、计算什么、提供什么结果。数学式采用可直接对照代码的记号。\par
\par\vspace{5pt}{\color{teal}\large 系统工作方式\par}
\[\begin{gathered}\text{场景 JSON}\to\text{类型校验}\to\text{几何诊断}\to\text{候选生成}\\ \to\text{独立复查}\to\text{选择改善方案}\to\text{下一轮或结束}\end{gathered}\]
几何校验器负责判定规则是否满足；大模型读取场景和非通过诊断，提出 MOVE 候选。当前模型请求只有文本 JSON，没有场景图片。规则策略也可以直接使用诊断内的 suggestions。\par
\par\vspace{5pt}{\color{teal}\large 阅读导航\par}
第 2 页：代码目录与数据模型\par 第 3 页：OBB、四元数、AABB 与轴对齐\par 第 4 页：SAT 碰撞与交叠体积\par 第 5 页：地面、接触、支撑面积与重心\par 第 6 页：支撑链、报告和参数\par 第 7 页：修复搜索、案例、接口与限制\par
\par\vspace{5pt}{\color{teal}\large 项目目录与版本\par}
{\footnotesize\color{ink} /Users/initial/Library/Mobile Documents/com\textasciitilde{}apple\textasciitilde{}CloudDocs/room\par}
{\footnotesize\color{ink} 本文中所有相对目录均以该路径为根目录。行号按生成时源码定位，后续编辑可能改变。\par}
{\footnotesize\color{ink} 分支对应仓库：github.com/initial3366-droid/room1\par 源码提交：b65f587\par}
适用边界：实心盒体、米制单位、统一世界坐标系、Z 轴向上。几何 pass 不等于真实物理稳定，不等于满足所有布局任务要求。\par
\clearpage
{\color{ink}\LARGE 01 目录索引与数据模型\par}\vspace{5pt}
{\footnotesize\color{ink} app/contracts/models.py\par}
基础合同：Box、CameraPose、VerificationResult；Pydantic 校验尺寸、四元数及数据格式。\par
{\footnotesize\color{ink} app/verification/models.py\par}
GeometryObject、SceneSnapshot、Diagnostic、DiagnosisReport、MovePrescription、RepairResult。\par
{\footnotesize\color{ink} app/environment/geometry.py\par}
底层向量和旋转工具：dot、cross、rotate、axes、corners。\par
{\footnotesize\color{ink} app/verification/geometry.py\par}
几何规则主体：bounds、aligned、penetration、DeterministicGeometryVerifier。\par
{\footnotesize\color{ink} app/verification/repair.py\par}
候选应用、移动限制、全场景复查及逐轮选择。\par
{\footnotesize\color{ink} app/verification/llm.py\par}
LLMRepairPolicy：把场景和诊断发送给模型，解析移动候选。\par
{\footnotesize\color{ink} app/verification/io.py\par}
读取 JSON；把简写 position\_m/size\_m 转换为完整 geometry 结构。\par
{\footnotesize\color{ink} app/verification/port.py\par}
verify\_layout：工件、版本和任务范围校验，应用布局后调用诊断。\par
{\footnotesize\color{ink} app/api/geometry.py\par}
HTTP 入口 /geometry/repair；调用修复流程并保存事件。\par
{\footnotesize\color{ink} configs/geometry-verifier-v1.json\par}
规则参数；Demo 场景见同目录 geometry-diagnosis-demo.json。\par
\par\vspace{5pt}{\color{teal}\large 核心数据流\par}
GeometryObject：object\_id、geometry、movable、anchored、requires\_support、support\_id、可选局部重心。SceneSnapshot：对象集合、frame\_id、floor\_z\_m、scene\_id。\par
Diagnostic：规则、状态、涉及对象、原因、测量值、严重程度、可编辑变量和建议。DiagnosisReport：场景修订标识、验证器版本、诊断集合和总体状态。\par
MovePrescription：object\_id + delta\_m + diagnosis\_id。RepairResult：最终场景、报告、接受的动作、迭代计数和结束状态。\par
{\footnotesize\color{ink} 输入合同拒绝重复对象 ID、自支撑、非法尺寸及坐标系不一致；support\_id 缺省允许保留为 None，由诊断返回 unknown。\par}
\clearpage
{\color{ink}\LARGE 02 盒体、旋转与边界\par}\vspace{5pt}
{\footnotesize\color{ink} 代码位置：app/contracts/models.py:36 · Box\par}
{\footnotesize\color{ink} 代码位置：app/verification/geometry.py:27 · bounds\par}
\par\vspace{5pt}{\color{teal}\large OBB：有方向的三维盒体\par}
中心 c = (cx, cy, cz)，尺寸 s = (sx, sy, sz)，半尺寸 h = s/2。旋转矩阵 R 的三列 u1、u2、u3 是局部单位轴。\par
\[\begin{gathered}\mathbf p=\mathbf c+\sum_{i=1}^{3}\alpha_i h_i\mathbf u_i,\qquad -1\le\alpha_i\le1\\ \text{八个角点：}\quad \alpha_i\in\{-1,+1\}\end{gathered}\]
position\_m 表示中心，size\_m 表示完整边长。一个边长 1 米、中心 (0,0,0.5) 的无旋转盒体，其 Z 范围为 [0,1]。\par
\par\vspace{5pt}{\color{teal}\large 四元数与旋转\par}
orientation\_xyzw 存储 q=(x,y,z,w)，合同要求其模长平方接近 1。对局部向量 v，旋转后得到 Rv；局部重心也通过相同旋转转换到世界坐标。\par
\[\begin{gathered}x^2+y^2+z^2+w^2\approx1\\ R=\begin{bmatrix}1-2(y^2+z^2)&2(xy-zw)&2(xz+yw)\\2(xy+zw)&1-2(x^2+z^2)&2(yz-xw)\\2(xz-yw)&2(yz+xw)&1-2(x^2+y^2)\end{bmatrix}\end{gathered}\]
\par\vspace{5pt}{\color{teal}\large AABB：世界坐标轴上的范围\par}
\[\begin{gathered}l_j=\min_{1\le k\le8}p_{k,j},\qquad u_j=\max_{1\le k\le8}p_{k,j},\quad j\in\{x,y,z\}\end{gathered}\]
bounds() 返回 (low, high)。旋转盒体的 AABB 会包含额外空间，不能直接用其重叠代替精确 OBB 碰撞检测。最低 Z 坐标用于地面检测。\par
{\footnotesize\color{ink} 代码位置：app/verification/geometry.py:44 · aligned\par}
\[\begin{gathered}\forall i\in\{1,2,3\}:\quad\max\{|u_{i,x}|,|u_{i,y}|,|u_{i,z}|\}\ge1-\varepsilon_{\mathrm{axis}}\end{gathered}\]
三轴均满足时视为轴对齐。90° 轴置换也可满足；普通 45° 旋转不满足。当前支撑规则只在轴对齐条件下给出确定结果。\par
\clearpage
{\color{ink}\LARGE 03 SAT 碰撞检测\par}\vspace{5pt}
{\footnotesize\color{ink} 代码位置：app/verification/geometry.py:61 · penetration\par}
\par\vspace{5pt}{\color{teal}\large 功能与候选分离轴\par}
分离轴定理：如果存在一个方向，使两个凸盒体的投影分离，那么盒体不穿透。代码测试 A 的 3 个轴、B 的 3 个轴，以及两组轴两两叉乘产生的 9 个方向，最多 15 个。长度小于 1e-10 的方向跳过，其余归一化为 n。\par
\[\begin{gathered}\mathcal N=\{\mathbf u_i\}\cup\{\mathbf v_j\}\cup\{\mathbf u_i\times\mathbf v_j\}\\ r_A(\mathbf n)=\sum_{i=1}^{3}h_{A,i}|\mathbf n\cdot\mathbf u_i|,\quad r_B(\mathbf n)=\sum_{i=1}^{3}h_{B,i}|\mathbf n\cdot\mathbf v_i|\\ d=(\mathbf c_B-\mathbf c_A)\cdot\mathbf n,\qquad p=r_A(\mathbf n)+r_B(\mathbf n)-|d|\end{gathered}\]
rA、rB 是投影半径，d 是中心在轴上的有符号距离。任意轴 p \(\le\) \(\varepsilon\)penetration 时返回 None；所有轴均超过容差才判定穿透。这里 p 表示沿该轴达到分离所需的移动量，在包含情形下不应混同于投影交集长度。\par
\par\vspace{5pt}{\color{teal}\large 分离移动候选\par}
\[\begin{gathered}\Delta\mathbf c_A=\begin{cases}-p\mathbf n,&d\ge0,\\+p\mathbf n,&d<0,\end{cases}\qquad \Delta\mathbf c_B=-\Delta\mathbf c_A\end{gathered}\]
返回按深度排序的 (depth, delta) 列表。建议只对当前物体对有效；必须复查它是否让对象撞到其他物体，或破坏原有支撑。\par
\par\vspace{5pt}{\color{teal}\large 例子：两把椅子重叠\par}
两个无旋转 1 米盒体的中心 X 坐标为 0 和 0.8，投影半径均为 0.5。p = 0.5 + 0.5 - 0.8 = 0.2 米。向左移 A 0.2 米，或向右移 B 0.2 米，可使这一对盒体刚好接触。\par
\par\vspace{5pt}{\color{teal}\large 轴对齐盒体交叠体积\par}
\[\begin{gathered}\ell_j=\max\{0,\min(u_{A,j},u_{B,j})-\max(l_{A,j},l_{B,j})\}\\ V=\ell_x\ell_y\ell_z\end{gathered}\]
仅双方均轴对齐时报告精确盒体交叠体积。一般旋转盒体体积为 None，SAT 仍能提供碰撞判断。全部物体对逐一检查，N 个物体有 N(N-1)/2 对，主要复杂度为 O(N\( {}^2 \))。\par
\clearpage
{\color{ink}\LARGE 04 地面与支撑几何\par}\vspace{5pt}
{\footnotesize\color{ink} 代码位置：app/verification/geometry.py:118 · diagnose\par}
\par\vspace{5pt}{\color{teal}\large 地面穿透\par}
\[\begin{gathered}d_{\mathrm{floor}}=z_{\mathrm{floor}}-l_z,\qquad d_{\mathrm{floor}}>\varepsilon_{\mathrm{penetration}}\Rightarrow\mathrm{fail}\\ \Delta\mathbf c=(0,0,d_{\mathrm{floor}})\end{gathered}\]
物体最低点 -0.1 米、地面 0 米，则向上移动 0.1 米。物体悬空仍可通过地面穿透项，需要另做支撑检查。\par
\par\vspace{5pt}{\color{teal}\large 垂直接触与水平重叠\par}
支撑顶面 zs 为地面高度，或声明支撑盒体的最高 Z。以下 L、U 是物体与支撑物在 XY 上的交集边界。\par
\[\begin{gathered}g=l_z-z_s\\ L_j=\max(l_{\mathrm{object},j},l_{\mathrm{support},j}),\quad U_j=\min(u_{\mathrm{object},j},u_{\mathrm{support},j})\\ A=\max(0,U_x-L_x)\max(0,U_y-L_y)\\ \mathrm{contact}=(|g|\le\varepsilon_{\mathrm{contact}})\land(A>0)\end{gathered}\]
g 为正表示悬空，为负表示进入支撑物。接触不仅要求高度接近，也要求水平方向有重叠。地面按无限平面处理，其投影面积取物体底面范围。\par
\par\vspace{5pt}{\color{teal}\large 重心与支撑边距\par}
\[\begin{gathered}\mathbf m=\mathbf c+R\mathbf m_{\mathrm{local}},\qquad\mathbf m_{\mathrm{local}}=\mathbf0\ \text{（未提供时）}\\ M=\min\{m_x-L_x,U_x-m_x,m_y-L_y,U_y-m_y\}\\ \mathrm{valid}=\mathrm{contact}\land(M\ge M_{\min})\end{gathered}\]
M<0 表示重心投影越界，M=0 表示位于边缘。地面支撑不计算有限边距；输入未给重心时假设均匀实心盒体。此判定是几何必要条件，不包含外力或动力学。\par
\par\vspace{5pt}{\color{teal}\large 修复建议与边界\par}
\[\begin{gathered}\Delta z=-g\\ \Delta x=\frac{l_{\mathrm{support},x}+u_{\mathrm{support},x}}{2}-m_x,\qquad \Delta y=\frac{l_{\mathrm{support},y}+u_{\mathrm{support},y}}{2}-m_y\end{gathered}\]
水平修正将重心移向支撑盒体中心，不保证最短。缺少 support\_id 返回 unknown；指定的支撑对象不存在返回 fail；非轴对齐支撑返回 unknown。anchored 或 requires\_support=false 的对象跳过局部支撑项。\par
\clearpage
{\color{ink}\LARGE 05 支撑链、报告与参数\par}\vspace{5pt}
{\footnotesize\color{ink} 代码位置：app/verification/geometry.py:223 · grounded\par}
\par\vspace{5pt}{\color{teal}\large 支撑关系是有向图\par}
\[\begin{gathered}G=(V,E),\qquad A\to B:\ \text{A 由 B 支撑}\\ \text{花瓶}\to\text{桌子}\to\mathrm{floor}\end{gathered}\]
递归到 floor 或声明的锚定对象返回 pass。节点不存在、循环或缺少可验证局部状态返回 unknown；遇到局部失败则传播其状态。seen 集合避免循环递归。代码只为局部支撑通过的对象追加 support\_chain 诊断。\par
例如上层箱子接触桌面，但桌子悬空：箱子的局部 support 可以 pass，沿桌子向下追溯的支撑链会 fail。锚定是输入声明，不自动验证真实锚固结构。\par
{\footnotesize\color{ink} 代码位置：app/verification/geometry.py:137 · add\par}
\par\vspace{5pt}{\color{teal}\large 报告与严重程度\par}
\[\begin{gathered}s_i=\min\left(1,\frac{\max(0,a_i)}{\lambda}\right)\\ \mathrm{status}=\begin{cases}\mathrm{fail},&\exists i:\ s_i^{\mathrm{status}}=\mathrm{fail},\\\mathrm{unknown},&\text{无 fail，但存在 unknown},\\\mathrm{pass},&\text{其余情形}.\end{cases}\end{gathered}\]
add() 统一生成 diagnosis\_id、measurements、editable\_variables 和 suggestions。碰撞以穿透深度作为 amount，地面以穿地深度作为 amount，支撑以间隙或边距不足量计算。严重程度是排序指标，不是事故概率。\par
\par\vspace{5pt}{\color{teal}\large 当前默认配置\par}
{\footnotesize\color{ink} 目录：configs/geometry-verifier-v1.json\par}
penetration\_tolerance\_m = 0.000001 m：小于或等于容差的穿透按接触处理。\par
contact\_tolerance\_m = 0.001 m：垂直间隙允许误差。\par
minimum\_support\_margin\_m = 0.001 m：重心离接触边缘的最小距离。\par
axis\_tolerance = 1e-9：轴对齐判定的无量纲容差。\par
severity\_scale\_m = 0.1 m：严重程度达到 1 的归一化尺度。\par
maximum\_move\_m = 2.0 m：每次候选移动距离上限，不是总行程预算。\par
maximum\_iterations = 20：默认循环上限；HTTP 请求默认 5，可指定 1 至 20。\par
{\footnotesize\color{ink} load\_config() 读取并校验配置；验证器 version 含配置指纹，scene.revision 含场景内容指纹，便于追溯和防止旧诊断作用于新场景。\par}
\clearpage
{\color{ink}\LARGE 06 修复搜索与使用说明\par}\vspace{5pt}
{\footnotesize\color{ink} 代码位置：app/verification/repair.py:41 · repair\_scene\par}
{\footnotesize\color{ink} 代码位置：app/verification/repair.py:17 · apply\_move\par}
\par\vspace{5pt}{\color{teal}\large 受约束的逐轮局部搜索\par}
\[\begin{gathered}J(S)=\left(N_{\mathrm{fail}},N_{\mathrm{unknown}},\sum_{i:\,\mathrm{status}_i=\mathrm{fail}}s_i\right)\\ J(S_{\mathrm{new}})<_{\mathrm{lex}}J(S_{\mathrm{old}}),\qquad \mathcal B(S_{\mathrm{new}})\subseteq\mathcal B(S_{\mathrm{old}})\\ \|\Delta\mathbf c\|_2\le D_{\max}\end{gathered}\]
按字典序先减少失败数量，再减少未知数量，最后比较失败严重程度。候选必须针对当前失败诊断的可编辑对象，固定物体不能移动；只更新位置，保留尺寸与旋转。每个候选独立应用于同一个当前场景，再全场景复查。\par
多个合格候选先比较 J，再比较移动长度，最后用移动 JSON 文本打破平局。每轮接受一个候选。全部规则通过返回 pass；无改善候选返回 blocked；耗尽轮数仍未通过返回 iteration\_limit。它不保证全局最优，也不会尝试先变差再变好的路径。\par
\par\vspace{5pt}{\color{teal}\large 模型如何知道哪里出错\par}
{\footnotesize\color{ink} 代码位置：app/verification/llm.py:68 · propose\par}
模型上下文包含 scene、非 pass 的 diagnostics、最大移动量和候选数量。诊断已给出错误对象、原因、测量值及规则建议；模型返回 moves。最终几何是否通过由 diagnose() 决定。\par
\par\vspace{5pt}{\color{teal}\large Demo 对应结果\par}
configs/geometry-diagnosis-demo.json：两把椅子重叠 0.2 米；floating\_box 底面距地面 0.75 米；vase 的重心超出 pedestal 支撑区域约 0.1 米。相应规则建议包括椅子横移 0.2 米、箱子下降 0.75 米、花瓶向底座中心横移 0.6 米。具体接受顺序由候选和复查结果决定。\par
\par\vspace{5pt}{\color{teal}\large 接口与复现入口\par}
app/api/geometry.py：POST /geometry/repair 接收 scene 与 max\_iterations。app/verification/port.py：verify\_layout 校验工件哈希、版本、任务范围，应用候选布局后诊断；额外任务约束仍报告 unknown。\par
{\footnotesize\color{ink} tests/test\_geometry\_verifier.py 覆盖碰撞、包含、支撑、修复和回放；tests/test\_llm\_repair.py 检查模型协议与移动门禁。现成调用脚本为 scripts/demo\_llm\_repair.py。\par}
\par\vspace{5pt}{\color{teal}\large 结论的准确范围\par}
pass 仅表示当前盒体规则通过。系统未计算摩擦、受力、质量分布演化、变形和移动路径碰撞；未识别真实图像，也未评价美观、人体工学或任意任务语义。文档描述当前源码实现，不构成真实设备的执行许可。\par
\end{document}